\documentclass[tikz,border=10pt]{standalone}
\usepackage{amsmath}
\usepackage{amssymb}
\usepackage{tikz}
\usetikzlibrary{shapes.geometric, arrows.meta, positioning, fit, backgrounds, shadows, patterns, decorations.pathmorphing, calc}

\begin{document}

\tikzset{
    box/.style={
        rectangle,
        rounded corners=3mm,
        draw=#1!70,
        fill=#1!15,
        thick,
        minimum height=1.2cm,
        text width=3.2cm,
        align=center,
        font=\small\bfseries,
        drop shadow
    },
    layer/.style={
        rectangle,
        rounded corners=2mm,
        draw=blue!60,
        fill=blue!8,
        minimum height=3cm,
        minimum width=0.8cm,
        font=\tiny
    },
    neuron/.style={
        circle,
        draw=orange!70,
        fill=orange!20,
        minimum size=0.4cm
    },
    arrow/.style={
        -Stealth,
        very thick,
        draw=gray!70
    },
    physarrow/.style={
        -Stealth,
        very thick,
        draw=red!70,
        dashed
    },
    label/.style={
        font=\scriptsize\bfseries,
        text=blue!70!black
    },
    math/.style={
        font=\footnotesize,
        text=black
    },
    physics/.style={
        rectangle,
        rounded corners=2mm,
        draw=red!70,
        fill=red!10,
        thick,
        minimum height=1.8cm,
        text width=4cm,
        align=center,
        font=\scriptsize\bfseries,
        drop shadow
    }
}

\begin{tikzpicture}[node distance=1.5cm and 1cm]

    % Input parameters
    \node[box=green, minimum height=3cm] (input) {
        Input\\Parameters\\
        $\mathbf{y}_n \in \mathbb{R}^P$
    };

    \node[font=\tiny, below=0.1cm of input, text width=3cm, align=center] (params) {
        $D, K, J, H_{\text{ext}}, T$
    };

    % Neural Network Architecture
    \node[layer, right=of input, minimum height=3.5cm, minimum width=1.2cm] (dense1) {};
    \node[font=\tiny, rotate=90, anchor=south] at (dense1.center) {Dense Layers};

    % Feature map
    \node[box=purple, right=of dense1, minimum height=2.5cm] (features) {
        Feature\\Embedding\\
        $\mathbf{F}_0$
    };

    % Transposed Convolutions
    \node[layer, right=0.8cm of features, minimum height=3cm, minimum width=1cm] (conv1) {};
    \node[font=\tiny, rotate=90, anchor=south] at (conv1.center) {Conv$^T$ Block 1};

    \node[layer, right=0.4cm of conv1, minimum height=2.5cm, minimum width=1cm] (conv2) {};
    \node[font=\tiny, rotate=90, anchor=south] at (conv2.center) {Conv$^T$ Block 2};

    \node[layer, right=0.4cm of conv2, minimum height=2cm, minimum width=1cm] (conv3) {};
    \node[font=\tiny, rotate=90, anchor=south] at (conv3.center) {Conv$^T$ Block 3};

    % Output image
    \node[box=orange, right=1cm of conv3, minimum height=3cm, minimum width=2.5cm] (output) {
        Predicted\\Spin Config.\\
        $\mathbf{\hat{I}}_n$
    };

    \node[font=\tiny, below=0.1cm of output, text width=2.3cm, align=center] (imgsize) {
        $H \times W \times C$
    };

    % LLG Physics Constraint Box (positioned between network and output)
    \node[physics, above=1.2cm of conv2] (llg) {
        LLG Equation\\
        $\frac{\partial \mathbf{m}}{\partial t} = -\gamma (\mathbf{m} \times \mathbf{H}_{\text{eff}})$\\
        $- \alpha \mathbf{m} \times \frac{\partial \mathbf{m}}{\partial t}$
    };

    % Collocation points
    \node[box=cyan, above=0.3cm of llg, minimum width=4cm, minimum height=0.8cm, text width=3.8cm] (colloc) {
        Collocation Points $\{t_i, \mathbf{r}_i\}$
    };

    % Draw arrows from physics to network layers
    \draw[physarrow] (llg.south) -- (conv2.north);
    \draw[physarrow] (colloc.south) -- (llg.north);

    % Forward pass arrows
    \draw[arrow] (input) -- (dense1);
    \draw[arrow] (dense1) -- (features);
    \draw[arrow] (features) -- (conv1);
    \draw[arrow] (conv1) -- (conv2);
    \draw[arrow] (conv2) -- (conv3);
    \draw[arrow] (conv3) -- (output);

    % Mathematical expressions
    \node[math, above=0.2cm of features] (eq1) {$f_\Theta$};

    % Title
    \node[font=\Large\bfseries, text=blue!80!black, above=0.8cm of input] (title) {Physics-Informed Neural Network (PINN) Architecture};

    % Training objective box with three loss components
    \node[box=red, below=1.8cm of conv2, minimum width=8cm, minimum height=2cm, text width=7.5cm] (loss) {
        Training: $\Theta^* = \arg\min_\Theta \left[\mathcal{L}_{\text{data}} + \lambda_p \mathcal{L}_{\text{physics}} + \lambda_m \mathcal{L}_{\text{mag}}\right]$\\[0.15cm]
        \textit{Data loss:} $\mathcal{L}_{\text{data}} = \|\mathbf{I}_n - \mathbf{\hat{I}}_n\|^2$\\
        \textit{Physics loss:} $\mathcal{L}_{\text{physics}} = \|\frac{\partial \mathbf{m}}{\partial t} + \text{LLG}[\mathbf{m}]\|^2$\\
        \textit{Magnitude:} $\mathcal{L}_{\text{mag}} = \|\|\mathbf{m}\| - 1\|^2$
    };

    % Advantages/Limitations boxes
    \node[box=green, below=0.5cm of loss, minimum width=3.8cm, text width=3.5cm, minimum height=1.5cm] (adv) {
        \textbf{Strengths:}\\
        • Physical constraints enforced\\
        • Reduced non-physical solutions\\
        • Better generalization
    };

    \node[box=red, right=0.3cm of adv, minimum width=3.8cm, text width=3.5cm, minimum height=1.5cm] (lim) {
        \textbf{Limitations:}\\
        • Complex hyperparameter tuning\\
        • Higher computational cost\\
        • Derivative approximation errors
    };

\end{tikzpicture}

\end{document}
